\DocumentMetadata{
  pdfversion=2.0,
  pdfstandard=ua-2,
  lang=en-US,
  tagging=on,
  tagging-setup={math/setup=mathml-SE,tabsorder=structure}
}
\documentclass[11pt,letterpaper]{article}
\usepackage[margin=0.68in,includefoot]{geometry}
\usepackage{newcomputermodern}
\usepackage{amsmath,amscd,mathtools}
\providecommand{\square}{\mdlgwhtsquare}
\usepackage[hidelinks]{hyperref}
\hypersetup{
  pdftitle={Numerical Linear Algebra First-Year Exam, May 2023},
  pdfauthor={University of Florida Department of Mathematics},
  pdfsubject={Graduate mathematics examination},
  pdfdisplaydoctitle=true,
  bookmarksopen=true
}
\setcounter{secnumdepth}{0}
\setlength{\parindent}{0pt}
\setlength{\parskip}{0.28em}
\setlength{\emergencystretch}{3em}
\setlength{\leftmargini}{2.15em}
\setlength{\leftmarginii}{2.65em}
\setlength{\leftmarginiii}{3em}
\renewcommand{\labelenumi}{\textbf{\theenumi.}}
\pagestyle{empty}
\raggedbottom
\allowdisplaybreaks
\newenvironment{examproblems}[1]{%
  \begin{enumerate}%
  \setcounter{enumi}{#1}%
  \setlength{\topsep}{0.35em}%
  \setlength{\partopsep}{0pt}%
  \setlength{\itemsep}{0.48em}%
  \setlength{\parsep}{0.18em}%
}{\end{enumerate}}
\newenvironment{examsubparts}{%
  \begin{enumerate}%
  \setlength{\topsep}{0.18em}%
  \setlength{\partopsep}{0pt}%
  \setlength{\itemsep}{0.16em}%
  \setlength{\parsep}{0.1em}%
}{\end{enumerate}}
\newenvironment{examinstructions}{%
  \par\begingroup\itshape
}{%
  \par\endgroup\vspace{0.18em}
}
\makeatletter
\renewcommand\section{\@startsection{section}{1}{0pt}%
  {-1.25ex plus -.25ex minus -.1ex}%
  {0.55ex plus .1ex}%
  {\normalfont\large\bfseries}}
\renewcommand{\@maketitle}{%
  \newpage\null\vskip -1.2em%
  {\footnotesize\bfseries
    \href{https://www.ufl.edu/}{University of Florida}, \href{https://math.ufl.edu/}{Department of Mathematics}\par}%
  \vskip 0.18em%
  \tagstructbegin{tag=Title}%
    {\LARGE\bfseries\href{https://gma.math.ufl.edu/exams/numerical-linear-algebra-first-year/}{Numerical Linear Algebra First-Year Exam}, May 2023\par}%
  \tagstructend%
  \vskip 0.35em%
  \tagmcbegin{artifact}%
    \hrule height 1pt%
  \tagmcend%
  \vskip 0.55em%
}
\makeatother
\title{Numerical Linear Algebra First-Year Exam, May 2023}
\author{}
\date{}
\begin{document}
\maketitle
\thispagestyle{empty}
\begin{examinstructions}
Do 4 (four) problems.
\end{examinstructions}
\begin{examproblems}{0}
\item
This problem has two parts.
\begin{examsubparts}
\item[\textnormal{(a)}]
If~\(P\) is a projector, prove that \(\operatorname{null}(P)\cap\operatorname{range}(P)=\{0\}\) and \(\operatorname{null}(P)=\operatorname{range}(I-P)\).
\item[\textnormal{(b)}]
Prove that~\(P\) is an orthogonal projector if and only if it is Hermitian.
\end{examsubparts}
\item
Let \(\lVert\cdot\rVert\) be a subordinate (induced) matrix norm.
\begin{examsubparts}
\item[\textnormal{(a)}]
If~\(E\) is \(n\times n\) with \(\lVert E\rVert<1\), then show \(I+E\) is nonsingular and 
\[
\lVert(I+E)^{-1}\rVert\leq\frac{1}{1-\lVert E\rVert}.
\]
\item[\textnormal{(b)}]
If~\(A\) is \(n\times n\) invertible and~\(E\) is \(n\times n\) with \(\lVert A^{-1}\rVert\lVert E\rVert<1\), then show \(A+E\) is nonsingular and 
\[
\lVert(A+E)^{-1}\rVert\leq\frac{\lVert A^{-1}\rVert}{1-\lVert A^{-1}\rVert\lVert E\rVert}.
\]
\end{examsubparts}
\item
Define the matrices~\(A\) and~\(B\) by 
\[
A=\begin{pmatrix}1&2&3\\2&4&6\\1&2&3\\2&4&6\end{pmatrix},
\qquad
B=\begin{pmatrix}1&1&0\\0&1&3\\0&1&3\\2&0&0\end{pmatrix}.
\]
\begin{examsubparts}
\item[\textnormal{(a)}]
Find a singular value decomposition of~\(A\).
\item[\textnormal{(b)}]
Find a QR decomposition of~\(B\).
\end{examsubparts}
\item
Assume \(A\in\mathbb{R}^{m\times m}\).
\begin{examsubparts}
\item[\textnormal{(a)}]
Prove that \(\langle x,y\rangle_A=x^*Ay\) is an inner product on \(\mathbb{R}^m\) if and only if~\(A\) is symmetric and positive definite.
\item[\textnormal{(b)}]
Assume that~\(A\) is symmetric and positive definite. If \(x_*\) is the solution to \(Ax=b\) and \(\{p_1,\cdots,p_m\}\) is an orthonormal basis for \(\mathbb{R}^m\) with respect to \(\langle\ ,\ \rangle_A\) and \(x_*=\sum c_i p_i\), give a formula for the \(c_i\).
\end{examsubparts}
\item
Let matrix \(A\in\mathbb{C}^{m\times n}\) with \(n<m\). Let a vector \(b\in\mathbb{C}^m\), and let~\(r\) denote the residual vector \(r=b-Ax\).
\begin{examsubparts}
\item[\textnormal{(a)}]
Show that~\(x\) solves the least-squares problem \(\min\lVert b-Ax\rVert_2\) if and only if \(r\in\operatorname{Null}(A^*)\).
\item[\textnormal{(b)}]
Suppose~\(A\) is full rank, and describe how to find the least-squares solution using the QR decomposition of~\(A\).
\end{examsubparts}
\end{examproblems}
\end{document}
